% ============================================================
\section{Discusión: la tercera luz en el modelado de las curvas}
\label{sec:tercera_luz}
% ============================================================

\subsection{El efecto en los parámetros del ajuste}
La sección~\ref{sec:dilucion} mostró que las curvas de \tx\ tienen entre 15 y 50~\% de flujo que no viene
de la binaria. En el modelado esa luz se llama \emph{tercera luz}: un flujo constante, que no se eclipsa,
y que entra a la apertura desde estrellas vecinas, la nebulosa o una tercera estrella ligada al sistema.
Si una fracción $c$ del flujo normalizado es tercera luz, el modelo observado es
\begin{equation}
F_{\rm obs}(t)=(1-c)\,\frac{F_{\rm bin}(t)}{F_{\rm bin,0}}+c,
\label{eq:tercera_luz}
\end{equation}
con $F_{\rm bin,0}$ el flujo de la binaria fuera de eclipse, de modo que $D_{\rm obs}=(1-c)\,D$.

Los ajustes con starry de este trabajo suponen $c=0$: las amplitudes de las dos estrellas suman uno con
priors estrechos. Para reproducir un eclipse diluido, el modelo compensa con la geometría, ya sea con una
secundaria más pequeña ($k=R_B/R_A$ menor) o con un eclipse más rasante (inclinación menor). Como
$R_2\propto\sqrt{D}$, una tercera luz de $c\approx0.3$ se traduce en un $R_2$ subestimado en $\sim20$~\%,
y con él la masa y la temperatura de la secundaria que se obtienen de las isócronas. Además, como $c$
cambia de un sector a otro, un modelo con $c=0$ convierte esas diferencias en geometrías distintas para cada
sector, como se vio en TIC~427353034 S98 y TIC~50897998 S98.

\subsection{Cómo se ha abordado en otros trabajos}
Los enfoques publicados se pueden agrupar en cuatro, y en la práctica suelen combinarse.

\paragraph{Tercera luz como parámetro libre, por sector.} Es el tratamiento más extendido con TESS.
En la serie \emph{Rediscussion of eclipsing binaries}, Southworth ajusta la tercera luz con JKTEBOP como
un parámetro más, por separado en cada sector o periodo de observación \citep[por ejemplo,][]{Southworth2024CWEri}.
En el ajuste de tránsitos con \texttt{juliet} \citep{Espinoza2019}, la dilución es un parámetro por instrumento
y puede dejarse libre por sector, aunque muchos parámetros de dilución libres dificultan la convergencia.
PHOEBE~2.2 \citep{Jones2020} incorpora la tercera luz como flujo o como fracción del total.

\paragraph{Información externa para romper la degeneración.} La tercera luz está correlacionada con la
inclinación y con $k$, sobre todo en eclipses parciales, así que suele acotarse con información que no viene
de la curva de luz: el cociente de flujos espectroscópico \citep{David2016}; ajustes en varias bandas, porque
$c$ depende del color de la contaminante y la geometría no \citep[en Orión,][ajustan $\ell_3$ por banda y obtienen
$\ell_3=0.31$ en TESS para nuestro sistema~6]{KounkelStassun2024}; o la dinámica de una terciaria, a partir de los
tiempos de eclipse y la velocidad sistémica \citep{GomezMaqueoChew2019}.

\paragraph{Quitar la contaminación en la fotometría.} Con fotometría PSF que modela y resta las estrellas de
Gaia, como TGLC \citep{HanBrandt2023}, la dilución se corrige antes del ajuste; \citet{Han2025} muestran que, sin
esto, los radios planetarios de TESS salen subestimados en un 6~\% en promedio. SPOC corrige con la fracción de
contaminación del TIC (\texttt{CROWDSAP}), que \citet{Baum2025} revisan sector por sector antes de interpretar
cambios de profundidad. En ajustes conjuntos con fotometría terrestre de alta resolución angular, la dilución de
ésta se fija en uno y solo se ajusta la de TESS.

\paragraph{Comparación de modelos.} Varios trabajos ajustan con y sin tercera luz y eligen con un criterio
estadístico (por ejemplo, el BIC), de modo que solo se incluye cuando los datos lo justifican.

\subsection{Aplicación a esta muestra}
No contamos con cocientes de flujo espectroscópicos ni con fotometría terrestre, pero sí con dos elementos que
rara vez se tienen juntos: varios sectores con contaminación distinta y una medición independiente de $c$ por
sector, la comparación con TGLC (tabla~\ref{tab:dilucion}). El tratamiento más directo sería ajustar una sola
geometría para todos los sectores de un sistema, con un $c$ propio en cada sector y un prior centrado en el
valor de TGLC (con $\sigma_c\approx0.1$, para cubrir las diferencias entre versiones de TGLC), y comprobarlo
con el ajuste de las curvas TGLC de S98 con $c\approx0$. Este esquema tiene además una ventaja diagnóstica:
un $c$ por sector no puede reproducir un cambio de $D_2/D_1$, de modo que en los sistemas 1, 9, 12 y 18 el
cambio de $D_2/D_1$ debería seguir apareciendo como residuo si es real.

Mientras no se incluya la tercera luz en el modelo, los radios y masas de las secundarias derivados aquí deben
tomarse como \textbf{límites inferiores}.

% ============================================================
\section{Discusión: la luminosidad de la primaria del sistema 11 (TIC 138829403)}
\label{sec:sistema11}
% ============================================================
El tipo espectral de la literatura para la primaria del sistema~11 es B3.5 (FAST; \texttt{Macro\_21aug2026.csv}),
que en la escala de Pecaut \& Mamajek corresponde a $T_{\rm eff}\approx16\,700$~K. Con la fotometría de Gaia y
2MASS, la paralaje corregida ($\varpi=2.336\pm0.025$~mas, $d\approx430$~pc, RUWE $=1.34$) y una extinción de
$A_V=1.3$ (MassAge, anclado en $J$), la luminosidad resulta $\log L/L_\odot=2.25$. A 10~Myr, la isócrona PARSEC
de secuencia principal predice para esa temperatura $\log L/L_\odot=2.62$ ($M\approx4.6\,M_\odot$,
$R\approx2.5\,R_\odot$): la estrella queda $0.37$~dex (un factor 2.3) \textbf{por debajo} de la secuencia principal,
en una zona del diagrama HR que ninguna isócrona ocupa. Por eso MassAge no encuentra solución en la fase
pre-secuencia principal y, en la fase completa, devuelve una edad sin sentido ($\sim$100--330~Myr).

Las explicaciones posibles son:
\begin{itemize}
\item \textbf{Extinción subestimada.} Para subir la luminosidad 0.37~dex hacen falta $\sim0.9$~mag más en la banda
  $J$, es decir $A_V\approx3$--4 en lugar de 1.3. MassAge estima $A_V$ comparando colores observados con los
  intrínsecos de la temperatura dada; si la $T_{\rm eff}$ es menor que la supuesta, la extinción derivada también
  se subestima.
\item \textbf{Tipo espectral o temperatura sobreestimados.} El TIC~8.2 da $T_{\rm eff}=8\,388$~K a partir de colores,
  muy distinto del B3.5 espectroscópico. Si la estrella fuera más fría, con la misma fotometría quedaría en otra
  posición del diagrama HR. La $T_{\rm eff}$ del TIC no está corregida por enrojecimiento y no es confiable aquí,
  pero la discrepancia sugiere revisar la clasificación espectral.
\item \textbf{Distancia.} Para explicar el déficit con la distancia, la paralaje tendría que ser $\sim1.5$ veces
  menor; con $\sigma_\varpi/\varpi\approx1$~\% y RUWE de 1.34 eso es poco probable.
\end{itemize}

Este problema afecta a los parámetros de la primaria y, a través de ellos, a los de la secundaria: el radio
$R_1=1.63\,R_\odot$ del catálogo (Stefan--Boltzmann) es menor que los 2.0--2.2~$R_\odot$ del ajuste con starry y que
los $\sim2.5\,R_\odot$ esperados para una B3.5, de modo que $R_2=R_1\sqrt{D_1}$ y la masa de la secundaria que se
deriva de las isócronas ($\sim0.06\,M_\odot$) están subestimados. No afecta, en cambio, a la edad asignada al
sistema, que proviene de sus hermanas cinemáticas ($10.8\pm1.8$~Myr) y no de la posición de la primaria en el
diagrama HR. Resolverlo requiere un espectro que confirme el tipo espectral y una estimación independiente de la
extinción (por ejemplo, a partir de mapas 3D de polvo o de los colores de estrellas vecinas a la misma distancia).

